% !PRINTED 346
\weberpage{346}

%<<CONT
and in our case of the quadratic field this reduces to the requirement
that
\weq{2}{\theta^p \equiv \theta \ (\mathrm{mod}\ \mathfrak{p})}
must hold.

If we take $\Delta = -4$, then $\theta = i = \sqrt{-1}$, and the congruence
 (2) becomes
\[
i^{p-1} \equiv 1 \ (\mathrm{mod}\ \mathfrak{p}).
\]

Now since $(1+i)(1-i) = 2$, so that for an odd $p$ neither 2 nor
$(1 \pm i)$ can be divisible by $\mathfrak{p}$, we must have $i^{p-1} = 1$,
and hence $p - 1$ must be divisible by 4.

Accordingly there follows from \S\,92, 2.:

1) \gesperrt{The congruence}
\[
x^2 \equiv -4 \ (\mathrm{mod}\ 4p),
\]
\gesperrt{or, what is the same thing,}
\[
x^2 \equiv -1 \ (\mathrm{mod}\ p)
\]
\gesperrt{has solutions then and only then, when $p \equiv 1 \ (\mathrm{mod}\ 4)$.}

If we take, secondly, $\Delta = 8$, then $\theta = \sqrt{2}$, or, if
$r = \sqrt{i}$ is an 8th root of unity,
\[
\theta = r + r^{-1},
\]
and by the binomial theorem the condition (2) can be represented
thus:
\[
r^p + r^{-p} \equiv r + r^{-1} \ (\mathrm{mod}\ \mathfrak{p}).
\]

This requires:
\[
(r^p - r)(r^p - r^{-1})\, r^{-p} \equiv 0 \ (\mathrm{mod}\ \mathfrak{p}),
\]
hence:
\[
r^p \equiv r \ \text{or} \ r^p \equiv r^{-1} \ (\mathrm{mod}\ \mathfrak{p}),
\]
\[
r^{p \pm 1} \equiv 1, \quad i^{\frac{p \pm 1}{2}} - 1 \equiv 0 \ (\mathrm{mod}\ \mathfrak{p}),
\]
and from this it follows as above, for an odd $p$:
\[
p \equiv \pm 1 \ (\mathrm{mod}\ 8).
\]

2) \gesperrt{The congruence}
\[
x^2 \equiv 2 \ (\mathrm{mod}\ p)
\]
\gesperrt{is therefore solvable then and only then, when $p \equiv 1$ or
$p \equiv -1 \ (\mathrm{mod}\ 8)$.}

Finally we set, if $q$ is an odd prime,
\weq{3}{\Delta = \pm q}
%CONT>>
